\subsection{Ontology Matching}
\label{sec:ontology-matching}
\slidesrc{Ontology matching: 1--14; Ontology matching (long version): 1--112; Ontology
development: 75--76}

Different people build different ontologies of the same domain: they use different
names (\emph{PC} vs.\ \emph{Personal Computers}), different structures, different
languages. To make two systems that use two ontologies work together we must find
which entity of one corresponds to which entity of the other. This is
\textbf{ontology matching}. The source is the tutorial by Pavel Shvaiko and Jérôme
Euzenat (SWAP 2006); the slides shown in class are the definitions of correspondence and
alignment, the matching process and the application domains; the rest comes from the
long version.

\subsubsection{The matching problem}
\slidesrc{Ontology matching (long version): 2--25}

The goals of the tutorial: illustrate the role of ontology matching, give an overview of
the basic matching techniques, show how state-of-the-art systems use them, and motivate
future research.

\begin{nbdef}{Matching operation}
The \textbf{matching} operation takes as input ontologies, each consisting of a set of
discrete entities (tables, XML elements, classes, properties, \ldots), and determines as
output the \textbf{relationships} (equivalence, subsumption, \ldots) holding between
these entities.
\end{nbdef}

\begin{center}
\begin{tikzpicture}[n/.style={draw, rounded corners=2pt, font=\small, inner sep=2.5pt},
    a/.style={font=\small, text=nbgreen!60!black},
    c/.style={{Stealth}-{Stealth}, thick, nbblue},
    r/.style={font=\small, fill=white, inner sep=1pt}]
  % left schema
  \node[n, anchor=west] (e1)  at (0,0)     {Electronics};
  \node[n, anchor=west] (pc1) at (0.4,-0.8) {Personal Computers};
  \node[n, anchor=west] (mp1) at (0.8,-1.6) {Microprocessors};
  \node[a, anchor=west]       at (1.2,-2.3) {PID, Name, Quantity, Price};
  \node[n, anchor=west] (ac1) at (0.4,-3.0) {Accessories};
  \node[n, anchor=west] (ph1) at (0.4,-3.8) {Photo and Cameras};
  \node[a, anchor=west]       at (0.8,-4.5) {PID, Name, Quantity, Price};
  % right schema
  \node[n, anchor=east] (e2)  at (13,0)     {Electronics};
  \node[n, anchor=east] (pc2) at (12.6,-0.8) {PC};
  \node[n, anchor=east] (mp2) at (12.2,-1.6) {PC board};
  \node[a, anchor=east]       at (11.8,-2.3) {ID, Brand, Amount, Price};
  \node[n, anchor=east] (cp2) at (12.6,-3.0) {Cameras and Photo};
  \node[n, anchor=east] (ac2) at (12.2,-3.8) {Accessories};
  \node[n, anchor=east] (dc2) at (12.2,-4.6) {Digital Cameras};
  \node[a, anchor=east]       at (11.8,-5.3) {ID, Brand, Amount, Price};
  \draw[c] (e1) -- node[r] {$=$} (e2);
  \draw[c] (pc1) -- node[r] {$=$} (pc2);
  \draw[c] (mp1) -- node[r] {$\perp$} (mp2);
  \draw[c] (ph1.east) -- node[r, pos=0.5] {$=$} (cp2.west);
  \draw[c] (ph1.east) -- node[r, pos=0.6] {$\ge$} (dc2.west);
\end{tikzpicture}
\captionof{figure}{Motivation: two XML schemas of an electronics catalogue (redrawn,
simplified). Blue arrows are correspondences: equal ($=$), more general ($\ge$),
disjoint ($\perp$). Among the attributes (green): Name $\ge$ Brand, Quantity $=$
Amount, Price $=$ Price.}
\label{fig:xml-schemas}
\end{center}

\begin{nbexample}{reading Figure~\ref{fig:xml-schemas}}
\emph{Personal Computers} on the left is the same as \emph{PC} on the right (equivalence,
found e.g.\ by expanding the abbreviation). \emph{Microprocessors} and \emph{PC board}
are disjoint: nothing is both. \emph{Photo and Cameras} is more general than
\emph{Digital Cameras}. Among the attributes, \emph{Quantity} = \emph{Amount}
(synonyms), and \emph{Name} is more general than \emph{Brand} (a brand is a kind of
name). The second motivating example in the slides does the same with two
\emph{ontologies} of products and books (\emph{Product} vs.\ \emph{Monograph},
\emph{creator} vs.\ \emph{author}, \emph{Person} vs.\ \emph{Human}), including
instances such as ``Bertrand Russell: My life'' (an Autobiography) and ``Albert Camus:
La chute'' (Literature).
\end{nbexample}

\paragraph{Schema matching vs.\ ontology matching.}
\textbf{Differences}: schemas often do not give explicit semantics for their data;
relational schemas provide no generalisation (no is-a); ontologies are logical systems
that constrain the meaning, defined as sets of logical axioms. \textbf{Commonalities}:
both provide a vocabulary of terms that describes a domain of interest, and both
constrain the meaning of these terms. Techniques developed for both problems are of
mutual benefit.

\paragraph{Scope.}
Heterogeneity between ontologies can occur when different \emph{languages} are used
(OWL vs.\ XML schema), different \emph{terminologies} (\emph{car} vs.\
\emph{automobile}), different \emph{modelling} (colour as a property vs.\ as a class),
\ldots\ Reducing heterogeneity takes two steps: \textbf{match}, which determines the
alignment, and \textbf{process the alignment} (merging, transforming, \ldots). Matching
can happen at \textbf{design time} (by the developers, before the system runs) or at
\textbf{run time} (on the fly, e.g.\ when two agents meet for the first time).

\subsubsection{Correspondence, alignment, matching process}
\slidesrc{Ontology matching: 1--13; Ontology matching (long version): 26--38}

\begin{nbdef}{Correspondence}
Given two ontologies $O$ and $O'$, a \textbf{correspondence} $M$ between $O$ and $O'$ is
a 5-tuple $\langle id, e, e', R, n\rangle$ such that:
\begin{itemize}[nosep]
  \item $id$ is a unique \textbf{identifier} of the correspondence;
  \item $e$ and $e'$ are \textbf{entities} of $O$ and $O'$ (e.g.\ XML elements, classes);
  \item $R$ is a \textbf{relation}, e.g.\ equivalence ($=$), more general
        ($\sqsupseteq$), disjointness ($\perp$);
  \item $n$ is a \textbf{confidence measure} in some mathematical structure (typically
        in the $[0,1]$ range).
\end{itemize}
\end{nbdef}

\begin{nbdef}{Alignment}
Given two ontologies $O$ and $O'$, an \textbf{alignment} $A$ between $O$ and $O'$
\begin{itemize}[nosep]
  \item is a \textbf{set of correspondences} on $O$ and $O'$;
  \item with some \textbf{cardinality}: 1-1, 1-*, etc.;
  \item with some additional \textbf{metadata} (method, date, properties, etc.).
\end{itemize}
\end{nbdef}

\begin{nbexample}{correspondences from Figure~\ref{fig:xml-schemas}}
\begin{align*}
  &\langle c_1,\ \text{Personal Computers},\ \text{PC},\ =,\ 0.95\rangle\\
  &\langle c_2,\ \text{Microprocessors},\ \text{PC board},\ \perp,\ 0.80\rangle\\
  &\langle c_3,\ \text{Photo and Cameras},\ \text{Digital Cameras},\ \sqsupseteq,\ 0.70\rangle
\end{align*}
The set $\{c_1, c_2, c_3, \ldots\}$ is an alignment. A \textbf{1-1} alignment matches each
entity to at most one entity of the other ontology; a \textbf{1-*} alignment lets one
entity of $O$ match several entities of $O'$ (\emph{Photo and Cameras} matches both
\emph{Cameras and Photo} and \emph{Digital Cameras}). The confidence values here are
invented for the example.
\end{nbexample}

\begin{center}
\begin{tikzpicture}[o/.style={circle, fill=yellow!70, font=\small, minimum size=0.75cm,
    inner sep=0pt}, al/.style={circle, fill=nbgreen!50, font=\small, minimum size=0.75cm,
    inner sep=0pt}, arr/.style={-{Stealth}}]
  \node[o] (O) at (0,1.1) {$O$};
  \node[al] (A) at (0,0) {$A$};
  \node[o] (O2) at (0,-1.1) {$O'$};
  \node[draw, font=\small\itshape, minimum height=0.8cm] (m) at (2.3,0) {matching};
  \node[al] (A2) at (4.5,0) {$A'$};
  \node[font=\small\itshape] (p) at (2.3,1.2) {parameters};
  \node[font=\small\itshape] (r) at (2.3,-1.2) {resources};
  \draw[arr] (O) -- (m); \draw[arr] (A) -- (m); \draw[arr] (O2) -- (m);
  \draw[arr] (m) -- (A2); \draw[arr] (p) -- (m); \draw[arr] (r) -- (m);
\end{tikzpicture}
\captionof{figure}{The matching process.}
\label{fig:matching-process}
\end{center}

\begin{nbfact}{the matching process}
The matching process takes as input two ontologies $O$ and $O'$, optionally an
\emph{input alignment} $A$ (correspondences already known, to be completed), some
\textbf{parameters} (thresholds, weights) and external \textbf{resources} (a thesaurus
like WordNet, an upper ontology), and returns an alignment $A'$.
\end{nbfact}

\subsubsection{Application domains}
\slidesrc{Ontology matching: 14; Ontology matching (long version): 39--52}

\begin{itemize}[nosep]
  \item \textbf{Traditional}: ontology evolution (match the old and the new version of
        an ontology), schema integration, catalog integration, data integration.
  \item \textbf{Emergent} (long version): P2P information sharing, agent communication,
        Web service composition, query answering on the Web.
\end{itemize}

\begin{center}
\begin{tikzpicture}[o/.style={circle, fill=yellow!70, font=\small, minimum size=0.7cm,
    inner sep=0pt}, b/.style={draw, font=\small, minimum height=0.7cm},
    arr/.style={-{Stealth}}]
  \begin{scope}
  \node[o] (O) at (0,2.4) {$O$}; \node[o] (O2) at (4,2.4) {$O'$};
  \node[b] (M) at (2,2.4) {Matcher};
  \node[circle, fill=nbgreen!50, font=\small, inner sep=1pt] (A) at (2,1.5) {$A$};
  \node[b] (G) at (2,0.6) {Generator};
  \node[o] (DB) at (0,-0.4) {DB}; \node[o] (P) at (4,-0.4) {Portal};
  \node[b, fill=nbred!15, font=\small\itshape] (T) at (2,-0.4) {Transformation};
  \draw[arr] (O) -- (M); \draw[arr] (O2) -- (M); \draw[arr] (M) -- (A);
  \draw[arr] (A) -- (G); \draw[arr] (G) -- (T); \draw[arr] (DB) -- (T);
  \draw[arr] (T) -- (P); \draw[arr] (DB) -- (O); \draw[arr] (P) -- (O2);
  \node[font=\small\bfseries] at (2,-1.2) {catalog integration};
  \end{scope}
  \begin{scope}[xshift=6.3cm]
  \node[o] (O) at (0,2.4) {$O$}; \node[o] (O2) at (4,2.4) {$O'$};
  \node[b] (M) at (2,2.4) {Matcher};
  \node[circle, fill=nbgreen!50, font=\small, inner sep=1pt] (A) at (2,1.5) {$A$};
  \node[b] (G) at (2,0.6) {Generator};
  \node[b, fill=nborange!15] (p1) at (0,-0.4) {peer1};
  \node[b, fill=nborange!15] (p2) at (4,-0.4) {peer2};
  \node[b, fill=nbred!15, font=\small\itshape] (T) at (2,-0.4) {mediator};
  \draw[arr] (O) -- (M); \draw[arr] (O2) -- (M); \draw[arr] (M) -- (A);
  \draw[arr] (A) -- (G); \draw[arr] (G) -- (T);
  \draw[arr] (p1) -- (O); \draw[arr] (p2) -- (O2);
  \draw[arr] ([yshift=2pt]p1.east) -- node[above, font=\tiny] {query} ([yshift=2pt]T.west);
  \draw[arr] ([yshift=-2pt]T.west) -- node[below, font=\tiny] {answer} ([yshift=-2pt]p1.east);
  \draw[arr] ([yshift=2pt]T.east) -- node[above, font=\tiny] {query} ([yshift=2pt]p2.west);
  \draw[arr] ([yshift=-2pt]p2.west) -- node[below, font=\tiny] {answer} ([yshift=-2pt]T.east);
  \node[font=\small\bfseries] at (2,-1.2) {P2P information sharing};
  \end{scope}
\end{tikzpicture}
\captionof{figure}{Two applications (simplified). Left: the ontologies of a database and
of a portal are matched; from the alignment a generator produces the transformation that
moves data from one to the other. Right: two peers with different ontologies; the
alignment produces a mediator that translates queries and answers.}
\end{center}

\begin{center}
\small
\begin{tabular}{@{}l ccccc l@{}}
\toprule
\textbf{application} & \rotatebox{90}{instances} & \rotatebox{90}{run time} &
  \rotatebox{90}{automatic} & \rotatebox{90}{correct} & \rotatebox{90}{complete} &
  \textbf{operation} \\
\midrule
Ontology evolution & \checkmark & & & \checkmark & \checkmark & transformation \\
Schema integration & \checkmark & & & \checkmark & \checkmark & merging \\
Catalog integration & \checkmark & & & \checkmark & \checkmark & data translation \\
Data integration & \checkmark & & & \checkmark & \checkmark & query answering \\
P2P information sharing & & \checkmark & & & & query answering \\
Web service composition & & \checkmark & \checkmark & \checkmark & & data mediation \\
Multi agent communication & & \checkmark & \checkmark & \checkmark & \checkmark & data
  translation \\
Query answering & \checkmark & \checkmark & & & & query reformulation \\
\bottomrule
\end{tabular}
\captionof{table}{Requirements of each application (long version): does it need
instances, run-time matching, full automation, a correct alignment, a complete one; and
what is done with the alignment.}
\end{center}

\begin{nbintuition}[reading the table]
Traditional applications (top four) match at \emph{design time}: a human can check the
result, and the alignment must be \emph{correct and complete} because it is used once and
for all to transform or merge. Emergent applications match at \emph{run time}, often
\emph{automatically} (no human available when two agents meet): there, a partial
alignment is acceptable as long as what is found is correct.
\end{nbintuition}

\subsubsection{Classification of matching techniques}
\slidesrc{Ontology matching (long version): 53--65}

\textbf{Matching dimensions}:
\begin{itemize}[nosep]
  \item \emph{input}: the underlying models (XML, OWL, \ldots); \textbf{schema-level}
        (only classes and properties) vs.\ \textbf{instance-level} (also the data);
  \item \emph{process}: approximate vs.\ exact; how the input is interpreted;
  \item \emph{output}: cardinality (1-1, 1-*); only equivalence vs.\ diverse relations
        (subsumption, \ldots); graded (a confidence in $[0,1]$) vs.\ absolute confidence.
\end{itemize}

The classification of schema-based techniques has \textbf{three layers}: the
\emph{upper} layer is about the granularity of the match and how the input is
interpreted; the \emph{middle} layer lists the classes of elementary (basic) techniques;
the \emph{lower} layer classifies them by the kind of input they use.

\begin{center}
\footnotesize
\renewcommand{\arraystretch}{1.15}
\begin{tabularx}{\linewidth}{@{}>{\raggedright\arraybackslash}p{1.75cm}
  >{\raggedright\arraybackslash}p{1.45cm} >{\raggedright\arraybackslash}p{2.6cm}
  >{\raggedright\arraybackslash}X >{\raggedright\arraybackslash}p{2.0cm}@{}}
\toprule
\textbf{granularity} & \textbf{interpre\-tation} & \textbf{basic technique} &
  \textbf{what it uses} & \textbf{kind of input} \\
\midrule
\cellcolor{nbblue!8}Element-level & syntactic & string-based & name and description
  similarity & terminological \\
\cellcolor{nbblue!8} & syntactic & language-based & tokenisation, elimination &
  terminological (linguistic) \\
\cellcolor{nbblue!8} & syntactic & constraint-based & type similarity, key properties &
  structural (internal) \\
\cellcolor{nbblue!8} & external & linguistic resources & lexicons, thesauri (WordNet) &
  terminological (linguistic) \\
\cellcolor{nbblue!8} & external & upper / domain-specific formal ontologies & DOLCE,
  FMA & semantic \\
\midrule
\cellcolor{nborange!10}Structure-level & syntactic & graph-based & graph homomorphism,
  children, leaves & structural (relational) \\
\cellcolor{nborange!10} & syntactic & taxonomy-based & taxonomy structure & structural
  (relational) \\
\cellcolor{nborange!10} & external & repository of structures & structure metadata
  (previous matches) & structural \\
\cellcolor{nborange!10} & semantic & model-based & SAT solvers, DL reasoners &
  semantic \\
\bottomrule
\end{tabularx}
\captionof{table}{Classification of schema-based matching techniques (simplified),
rewritten as a table from the slide's diagram. Read each row from left to right: the
upper layer (granularity and interpretation), the middle layer (the basic technique), the
lower layer (the kind of input it works on). Upper/domain ontologies are used at both
granularities.}
\label{fig:matching-classification}
\end{center}

\begin{itemize}[nosep]
  \item \textbf{Element-level} techniques look at each entity in isolation (its name,
        its type); \textbf{structure-level} techniques look at how entities are related
        to the others (its children, its position in the taxonomy).
  \item \textbf{Syntactic} techniques use only the input itself; \textbf{external} ones
        use outside resources (WordNet, an upper ontology, a repository of previous
        matches); \textbf{semantic} ones reason on the formal meaning (logic).
\end{itemize}

\subsubsection{Basic techniques}
\slidesrc{Ontology matching (long version): 66--88}

\paragraph{Element-level, string-based.} (Used e.g.\ by COMA, SF, S-Match, OLA.)
\begin{itemize}[nosep]
  \item \textbf{Prefix}: checks whether the first string starts with the second:
        \emph{net} = \emph{network}; but also, wrongly, \emph{hot} = \emph{hotel}.
  \item \textbf{Suffix}: checks whether the first string ends with the second: \emph{ID}
        = \emph{PID}; but also, wrongly, \emph{word} = \emph{sword}.
  \item \textbf{Edit distance} (Section~\ref{sec:minimum-edit-distance}): the number of
        edit operations (insertions, deletions, substitutions of characters) needed to
        transform one string into the other, \textbf{normalised by the length of the
        longer string}.
\end{itemize}

\begin{nbexample}{EditDistance(NKN, Nikon) = 0.4}
Ignoring case, \emph{nkn} becomes \emph{nikon} with two insertions: n\,\underline{i}\,k\,%
\underline{o}\,n. So the edit distance is 2, and normalised by the length of the longer
string: $2/5 = 0.4$. A distance of 0.4 (similarity $1-0.4=0.6$) suggests that \emph{NKN}
may be an abbreviation of \emph{Nikon}. (Used e.g.\ by S-Match, OLA, Anchor-Prompt.)
\end{nbexample}

\paragraph{Element-level, language-based.} (COMA, Cupid, S-Match, OLA.)
\begin{itemize}[nosep]
  \item \textbf{Tokenisation}: parse names into tokens by recognising punctuation and
        case: \emph{Hands-Free Kits} $\to$ $\langle$hands, free, kits$\rangle$.
  \item \textbf{Lemmatisation}: morphological analysis of the tokens to find their basic
        forms: \emph{Kits} $\to$ \emph{Kit}.
  \item \textbf{Elimination}: discard ``empty'' tokens: articles, prepositions,
        conjunctions, \ldots (\emph{a, the, by, type of, their, from}), i.e.\ stop words
        (Section~\ref{sec:stop-words-tfidf}).
\end{itemize}

\paragraph{Element-level, linguistic resources.} Sense-based, with WordNet (Artemis,
CtxMatch, S-Match):
\begin{itemize}[nosep]
  \item $A \sqsubseteq B$ if $A$ is a hyponym or meronym of $B$: \emph{Brand}
        $\sqsubseteq$ \emph{Name};
  \item $A \sqsupseteq B$ if $A$ is a hypernym or holonym of $B$: \emph{Europe}
        $\sqsupseteq$ \emph{Greece};
  \item $A = B$ if they are synonyms: \emph{Quantity} $=$ \emph{Amount};
  \item $A \perp B$ if they are antonyms or siblings in the part-of hierarchy:
        \emph{Microprocessors} $\perp$ \emph{PC Board}.
\end{itemize}
\textbf{Gloss-based} (S-Match): compare the WordNet glosses of the two terms; the more
words they share, the higher the similarity, and equivalence is returned above a
threshold. E.g.\ \emph{Maltese dog}: ``a breed of toy dogs having a long straight silky
white coat''; \emph{Afghan hound}: ``a tall graceful breed of hound with a long silky
coat'' share \emph{breed, long, silky, coat}. This is the idea of the Lesk algorithm
(Section~\ref{sec:word-sense-disambiguation}).

\paragraph{Structure-level, taxonomy-based.} Ontologies are seen as graphs of terms and
their relations (Anchor-Prompt, NOM, QOM):
\begin{itemize}[nosep]
  \item \textbf{bounded path matching}: take two paths of hierarchical links, compare the
        terms and their positions along the paths, and identify similar terms;
  \item \textbf{super(sub)-concept rules}: if the super-concepts are the same, the
        concepts themselves are similar.
\end{itemize}

\paragraph{Structure-level, tree-based.} (Cupid, COMA.)
\begin{itemize}[nosep]
  \item \textbf{children}: two non-leaf elements are similar if their sets of immediate
        children are highly similar;
  \item \textbf{leaves}: two non-leaf elements are similar if their sets of leaves are
        highly similar, even if their immediate children are not.
\end{itemize}
E.g.\ \emph{Photos and cameras} (left) and \emph{Digital cameras} (right) both have the
leaves PID/ID, Name/Brand, Quantity/Amount, Price/Price: their leaves match, so the two
parents are probably related.

\paragraph{Structure-level, model-based.}
\textbf{Propositional satisfiability (SAT)} (CtxMatch, S-Match): the question ``are the
two nodes equivalent?'' is turned into a logical formula
$\text{Axioms} \to rel(context_1, context_2)$ and a SAT solver checks whether it is valid.
The \emph{context} of a node is the conjunction of the labels on its path from the root.

\begin{nbexample}{SAT-based matching of \emph{Personal Computers} and \emph{PC}}
Axioms from the lexicon: $\text{Electronics}_1 \leftrightarrow \text{Electronics}_2$ and
$\text{Personal Computers}_1 \leftrightarrow \text{PC}_2$. Contexts:
\[
  context_1 = \text{Electronics}_1 \wedge \text{Personal Computers}_1, \qquad
  context_2 = \text{Electronics}_2 \wedge \text{PC}_2.
\]
To prove equivalence we check that
\begin{align*}
  &(\text{Electronics}_1 \leftrightarrow \text{Electronics}_2) \wedge
   (\text{Personal Computers}_1 \leftrightarrow \text{PC}_2) \\
  &\qquad \to \bigl((\text{Electronics}_1 \wedge \text{Personal Computers}_1)
   \leftrightarrow (\text{Electronics}_2 \wedge \text{PC}_2)\bigr)
\end{align*}
is valid, i.e.\ that its negation is unsatisfiable. It is: if the two pairs of labels
are equivalent, the two conjunctions are equivalent. So the two nodes are equivalent.
\end{nbexample}

\textbf{Description logics (DL)-based}: use a DL reasoner on the definitions.

\begin{nbexample}{micro-company $\sqsubseteq$ SME}
\begin{align*}
  \text{micro-company} &= \text{company} \sqcap {\le_5}\,\text{employee}\\
  \text{SME} &= \text{firm} \sqcap {\le_{10}}\,\text{associate}
\end{align*}
Background axioms: company $=$ firm; associate $\sqsubseteq$ employee. A micro-company
has at most 5 employees; every associate is an employee, so it has at most 5 associates,
hence at most 10; and a company is a firm. Therefore every micro-company is an SME:
\[
  \text{micro-company} \sqsubseteq \text{SME}.
\]
No string or structure technique could find this: the names share nothing. Only
reasoning on the meaning does.
\end{nbexample}

\subsubsection{Composing matchers and selecting the alignment}
\slidesrc{Ontology matching (long version): 89--95}

Real systems combine several basic matchers.

\begin{center}
\begin{tikzpicture}[o/.style={circle, fill=yellow!70, font=\scriptsize, minimum size=0.6cm,
    inner sep=0pt}, al/.style={circle, fill=nbgreen!50, font=\scriptsize,
    minimum size=0.6cm, inner sep=0pt}, b/.style={draw, font=\scriptsize\itshape,
    minimum height=0.6cm}, arr/.style={-{Stealth}}]
  \node[o] (O) at (0,0.9) {$O$}; \node[al] (A) at (0,0) {$A$}; \node[o] (O2) at (0,-0.9) {$O'$};
  \node[b] (m1) at (1.6,0) {matching}; \node[al] (A1) at (3.0,0) {$A'$};
  \node[b] (m2) at (4.4,0) {matching$'$}; \node[al] (A2) at (5.8,0) {$A''$};
  \draw[arr] (O) -- (m1); \draw[arr] (A) -- (m1); \draw[arr] (O2) -- (m1);
  \draw[arr] (m1) -- (A1); \draw[arr] (A1) -- (m2); \draw[arr] (m2) -- (A2);
  \draw[arr] (O) to[bend left=25] (m2); \draw[arr] (O2) to[bend right=25] (m2);
  \node[font=\scriptsize\bfseries] at (2.9,-1.6) {sequential composition};
  \begin{scope}[xshift=7.4cm]
  \node[o] (O) at (0,0.9) {$O$}; \node[al] (A) at (0,0) {$A$}; \node[o] (O2) at (0,-0.9) {$O'$};
  \node[b] (m1) at (1.7,0.8) {matching}; \node[b] (m2) at (1.7,-0.8) {matching$'$};
  \node[al] (A1) at (3.0,0.8) {$A'$}; \node[al] (A2) at (3.0,-0.8) {$A''$};
  \node[b] (ag) at (4.3,0) {aggregation}; \node[al] (A3) at (5.6,0) {$A'''$};
  \foreach \s in {O,A,O2} {\draw[arr] (\s) -- (m1); \draw[arr] (\s) -- (m2);}
  \draw[arr] (m1) -- (A1); \draw[arr] (m2) -- (A2);
  \draw[arr] (A1) -- (ag); \draw[arr] (A2) -- (ag); \draw[arr] (ag) -- (A3);
  \node[font=\scriptsize\bfseries] at (2.8,-1.6) {parallel composition};
  \end{scope}
\end{tikzpicture}
\captionof{figure}{Composing matchers (parameters and resources of each matcher
omitted). Sequential: the output of the first matcher is the input alignment of the
second. Parallel: independent matchers whose results are aggregated.}
\end{center}

\textbf{Selecting the final alignment} from the candidate correspondences:
\begin{itemize}[nosep]
  \item \textbf{ranking strategies}: \emph{thresholds} (keep correspondences with
        confidence above $t$), \emph{MaxDelta} (keep those within $\delta$ of the best
        one);
  \item \textbf{cardinalities}: 1-1, 1-*, *-*;
  \item \textbf{optimisation}: \emph{stable marriage} (no two entities would both
        prefer each other to their assigned partners), \emph{maximal weight match}
        (maximise the total confidence);
  \item \textbf{directionality}: $O \to O'$ or $O' \to O$ (SmallLarge, LargeSmall), or
        both directions (Both).
\end{itemize}

\subsubsection{State-of-the-art systems}
\slidesrc{Ontology matching (long version): 96--104}

About 50 matching systems existed in 2006; four examples:

\begin{center}
\small
\begin{tabularx}{\linewidth}{@{}>{\raggedright\arraybackslash}p{2.6cm} >{\raggedright\arraybackslash}X >{\raggedright\arraybackslash}X@{}}
\toprule
\textbf{system} & \textbf{what it does} & \textbf{architecture} \\
\midrule
Cupid (U.\ Washington, Microsoft, U.\ Leipzig) & schema-based; similarity coefficients in
  $[0,1]$; linguistic and structure matching & sequential: linguistic matching (with
  thesauri) $\to$ structure matching $\to$ weighting $\to$ alignment \\
Falcon-AO (Southeast U., Nanjing; the slide says ``Southwest'') & linguistic and structure matching & linguistic
  matching $\to$ structure matching $\to$ alignment \\
OLA (INRIA Rhône-Alpes, U.\ de Montréal) & schema- and instance-based; computes
  dissimilarities, then extracts alignments (equivalences in $[0,1]$); terminological
  (incl.\ linguistic) and structural (internal and relational) distances & neither
  sequential nor parallel: the similarity computation is iterated until it stabilises \\
S-Match (U.\ Trento) & schema-based; computes $=$, $\sqsupseteq$, $\sqsubseteq$, $\perp$;
  analyses the meaning (concepts, not labels) codified in the elements and structures &
  sequential with composition at the element level: translator $\to$ pre-processing
  (with oracles) $\to$ PTrees $\to$ match manager (SAT solvers, semantic matchers) \\
\bottomrule
\end{tabularx}
\end{center}

\subsubsection{Conclusions, tools and benchmarks}
\slidesrc{Ontology matching (long version): 105--112; Ontology development: 75--76}

\textbf{Uses of the classification} of Table~\ref{fig:matching-classification}: it
gives a common basis to compare existing systems; it helps design a new system or
elementary matcher from the state of the art; it helps design systematic benchmarks
(e.g.\ discard features from ontologies one at a time and see which class of technique
deals with which feature).

\textbf{Challenges}: missing background knowledge; performance of the systems;
interactive approaches; explanations of matching; social aspects of ontology matching;
large-scale evaluation; infrastructures.

\textbf{Tools and benchmarks}: \url{http://www.ontologymatching.org} (the portal of the
field), a list of 50 ontology mapping and alignment tools, and the \textbf{OAEI}
(Ontology Alignment Evaluation Initiative, \url{http://oaei.ontologymatching.org/}), a
yearly campaign that evaluates matching systems on common benchmarks. The reference
book is Euzenat and Shvaiko, \emph{Ontology Matching} (Springer).

\begin{nbsummary}[ontology matching]
\begin{itemize}[nosep]
  \item Matching: from ontologies (sets of entities) to relationships between their
        entities. Heterogeneity of languages, terminologies, modelling. Match, then
        process the alignment; at design time or run time.
  \item Correspondence $\langle id, e, e', R, n\rangle$: identifier, two entities,
        relation ($=, \sqsupseteq, \perp$), confidence in $[0,1]$. Alignment: set of
        correspondences + cardinality + metadata.
  \item Matching process: $O, O'$, input alignment $A$, parameters, resources $\to A'$.
  \item Applications: ontology evolution, schema/catalog/data integration; P2P, agents,
        Web services, query answering.
  \item Techniques: element-level (string: prefix, suffix, edit distance; language:
        tokenisation, lemmatisation, elimination; resources: WordNet relations, glosses)
        and structure-level (taxonomy, tree: children/leaves, model: SAT, DL).
  \item Composition: sequential or parallel + aggregation; selection by thresholds,
        cardinality, optimisation, directionality. Systems: Cupid, Falcon-AO, OLA,
        S-Match. Benchmark: OAEI.
\end{itemize}
\end{nbsummary}
